\clearpage
\chapter*{Erratum à l'exposé n\textsuperscript{o} 149}
\addcontentsline{toc}{chapter}{Erratum à l'exposé n\textsuperscript{o} 149}
\setcounter{footnote}{0}
\src{298}

\noindent\textbf{1. Théorème de dualité pour les faisceaux algébriques
cohérents.}

\begin{center}
[Séminaire Bourbaki, t. 9, 1956/57, n\textsuperscript{o} 149, 25 p.]
\end{center}

\noindent\textbf{Remarque sur la page 14.} --- Comme je l'ai signalé dans ma
conférence au Congrès international des Mathématiciens en 1958,\footnote{
Grothendieck (Alexander). --- The cohomology theory of abstract algebraic
varieties, \emph{Proceedings of the International Congress of Mathematicians}
[1958, Edinburgh], p. 103--118. --- Cambridge, at the University Press, 1960.}
les questions soulevées ici sont maintenant complètement résolues.

Le lecteur trouvera d'autres renseignements sur la dualité des faisceaux
cohérents dans \emph{loco citato}\footnotemark[1], p. 112--115, dans EGA III,
2\textsuperscript{e} partie (en préparation), et, en attendant, dans SGA 1962.
Un traitement plus systématique se trouvera dans un chapitre ultérieur de EGA
(chapitre IX dans le plan prévu).
